\section{Finding Periods}
\label{sec:cses-finding-periods}

\problemstatement{A period of a string $s$ of length $n$ is an integer
$p$ ($1\le p\le n$) such that $s[i]=s[i+p]$ for all valid $i$ (equivalently,
$s$ is a prefix of an infinite repetition of its first $p$ characters).
Output all periods in increasing order.}

\begin{patternbox}
Periods and borders are two sides of one coin: \textbf{$p$ is a period of
$s$ if and only if $s$ has a border of length $n-p$}. So compute the
prefix function, enumerate all border lengths $b$ (including $0$), and
each gives a period $n-b$.
\end{patternbox}

\subsection*{The period--border duality}

If $s$ has a border of length $b$, the first $n-b$ characters repeat to
cover the string, making $n-b$ a period; conversely every period $p$
leaves a border of length $n-p$. The border lengths of the whole string
are the failure chain $\pi[n-1],\pi[\pi[n-1]-1],\ldots$ together with $0$
(the empty border, which corresponds to the trivial period $p=n$). Mapping
each border length $b$ to $n-b$ and sorting gives all periods.

\begin{ideabox}[Every Border Length $b$ Gives Period $n-b$]
Include the empty border ($b=0$) so that $p=n$ (the whole string is always
its own period) is present. The failure chain supplies the remaining
borders. One prefix-function pass yields every period.
\end{ideabox}

\subsection*{Algorithm}

\begin{enumerate}
  \item Compute the prefix function $\pi$.
  \item Collect border lengths via the failure chain from $\pi[n-1]$,
        plus $0$.
  \item For each border length $b$, output $n-b$; sort ascending.
\end{enumerate}

\subsection*{Dry run}

\dryrun{$s=$``abcabca'' ($n=7$).}

\begin{itemize}
  \item Longest border ``abca'' has length $4$; chain gives borders
        $4,1$; plus $0$.
  \item Periods: $7-4=3$, $7-1=6$, $7-0=7$.
  \item Increasing: $3,6,7$ (indeed ``abc'' repeats with period $3$).
\end{itemize}

\subsection*{C++ implementation}

\begin{lstlisting}
#include <bits/stdc++.h>
using namespace std;

vector<int> prefixFunction(const string& s) {
    int n = s.size();
    vector<int> pi(n, 0);
    for (int i = 1; i < n; i++) {
        int j = pi[i - 1];
        while (j > 0 && s[i] != s[j]) j = pi[j - 1];
        if (s[i] == s[j]) j++;
        pi[i] = j;
    }
    return pi;
}

int main() {
    string s;
    cin >> s;
    int n = s.size();
    vector<int> pi = prefixFunction(s);

    vector<int> periods;
    for (int b = pi[n - 1]; b > 0; b = pi[b - 1]) periods.push_back(n - b);
    periods.push_back(n);                         // border 0 -> period n
    sort(periods.begin(), periods.end());

    for (int p : periods) cout << p << " ";
    cout << "\n";
    return 0;
}
\end{lstlisting}

\begin{complexitybox}
\textbf{Time Complexity.} $O(n)$. \textbf{Space Complexity.} $O(n)$.
\end{complexitybox}

\begin{takeawaybox}
Period equals $n$ minus a border length. Once you see that borders and
periods are complementary, every ``how does this string repeat?'' question
becomes a prefix-function query. This duality also underlies detecting the
smallest repeating block of a string.
\end{takeawaybox}
